% Pacotes Essenciais ========================================================================================================
\usepackage{lipsum}
\usepackage{titlesec}

\usepackage[utf8]{inputenc} % ................................................................. Acentuação
\usepackage[brazilian]{babel} % .................................................................. Tradução para Pt-br
\usepackage[T1]{fontenc} % .................................................................... Fontes melhores
\usepackage{graphicx} 
\usepackage[svgnames,table]{xcolor}% ................................................................ Gráficos e Cores
\usepackage{subcaption}% permite sub figuras nas imagens

\usepackage[lmargin=2.5cm,tmargin=4cm,rmargin=2cm,bmargin=3.5cm,headheight=15.1pt]{geometry} % .... Margens da ABNT
\usepackage[onehalfspacing]{setspace} % ....................................................... Espaçamento de Um e Meio
\usepackage{indentfirst} % .................................................................... Indentar o primeiro
\usepackage{enumerate} % ...................................................................... Enumeração diferenciada
\usepackage{nameref} % ........................................................................ Referência por nome
\usepackage{anyfontsize} % .................................................................... Qualquer tamanho de fonte
\usepackage{blindtext,lipsum,comment} % ....................................................... Texto cego e comentário
\usepackage{multicol,multirow} % .............................................................. Multicolunas e multilinhas
\usepackage[hidelinks]{hyperref} % ....................................................................... Links
% ===========================================================================================================================



% TikZ ======================================================================================================================
\usepackage{tikz} % ........................................................................... O completíssimo pacote TIKZ
\usetikzlibrary{calc,patterns,arrows} % ....................................................... Algumas bibliotecas do Tikz
\usepackage{pgfplots} % ....................................................................... Para usar axis no Tikz
\pgfplotsset{compat=1.17} % ................................................................... Faz o pgfplots compatível
\usepgfplotslibrary{fillbetween} % ............................................................ Preenche entre funções
% ===========================================================================================================================



% Embelezamento das páginas =================================================================================================
\usepackage{fancyhdr} % ....................................................................... Fancy é para deixar bonito
\usepackage{graphics}

\renewcommand{\headrulewidth}{0pt} % .......................................................... Tamanho da régua superior
\renewcommand{\footrulewidth}{0pt} % .......................................................... Tamanho da régua inferior

\lhead{} % ........................................................................... Topo a esquerda da página 
\chead{} % .................................................................................... Topo ao centro da página 
\rhead{} % ............................................................................ Topo a direita da página 
\lfoot{} % .......................................................................... Baixo a esquerda da página 

\cfoot{{\color{cinza}
\footnotesize{\fontfamily{phv}\selectfont
\textbf{
Brasília, \today. \\
ELETRONJUN – UNIVERSIDADE DE BRASÍLIA \\
EMPRESA JÚNIOR DE ENGENHARIA ELETRÔNICA\\}
Área Especial de Indústria – Projeção A| 72444-240 | Gama/DF \\
www.eletronjun.com.br | contato@eletronjun.com.br}}
} % .................................................................................... Baixo ao centro da página 
\rfoot{} % .................................................................................... Baixo a direita da página 


\pagestyle{fancy} % ........................................................................... Carrega as páginas bonitas


% Novas fontes ==============================================================================================================
\usepackage{calligra}
% ===========================================================================================================================



% Novas cores ===============================================================================================================
\definecolor{vermelho}{HTML}{FF0000}
\definecolor{verde}{HTML}{00FF00}
\definecolor{azul}{HTML}{0000FF}
\definecolor{amarelo}{HTML}{FFFF00}
\definecolor{rosa}{HTML}{FF00FF}
\definecolor{aqua}{HTML}{00FFFF}
\definecolor{resposta}{HTML}{AD65C7}
\definecolor{cinza}{HTML}{595959}
\definecolor{verdeEJ}{HTML}{92D050}
% ===========================================================================================================================



% Bibliografia ==============================================================================================================
\usepackage[alf,abnt-repeated-title-omit=yes,abnt-emphasize=bf,abnt-etal-list=0]{abntex2cite} % Citação de acordo com a ABNT
\citebrackets() % ............................................................................. Citação com parêntesis
% ===========================================================================================================================
%NÃO SEI PQ COLOCARAM DOIS PACOTES RELACIONADOS A BIBLIOGRAFIA ENTÃO COMENTEI UM
% ---
% BIBLIOGRAFIA
% Pacotes de citações para o padrão ABNT
% ---
% Esse pacote usa o biblatex-abnt
%\usepackage[style=ieee,
%   %style=abnt, % Estilo padrão ABNT
%	language=brazil, % Lingua pt_BR
%	%noslsn, % Oculta as abreviações [s.l], [s.n] e [s.l.: s.n.]
%	repeatfields, % Imprime os campos repetidos na bibliografia, em vez de ___________.
%	sorting=none,
%	]{biblatex}
%\addbibresource{Ref.bib} 


\usepackage{pifont}% pacote para símbolo estranho do comando \ding{118} para saber mais leia a documantação em https://mirrors.mit.edu/CTAN/macros/latex/required/psnfss/psnfss2e.pdf
\newcommand{\pra}{\ding{118}}

% Matemática ================================================================================================================
\usepackage{amsfonts,amssymb,amsthm,amsmath,dsfont,mathtools} % ............................... Fontes de Matemática
\everymath{\displaystyle} % ................................................................... Matemática com estilo


\newcommand{\NN}{\mathds{N}} % ................................................................ Conjunto dos números Naturais
\newcommand{\ZZ}{\mathds{Z}} % ................................................................ Conjunto dos números Inteiros
\newcommand{\QQ}{\mathds{Q}} % ................................................................ Conjunto dos números Racionais
\newcommand{\RR}{\mathds{R}} % ................................................................ Conjunto dos números Reais
\newcommand{\CC}{\mathds{C}} % ................................................................ Conjunto dos números Complexos
\newcommand{\PP}{\mathds{P}} % ................................................................ Conjunto dos números Primos
\newcommand{\MM}{\mathds{M}} % ................................................................ Conjunto de matrizes
% ===========================================================================================================================



% Novos ambientes ===========================================================================================================
%\newenvironment{citelong} % ................................................................... Citação longa
%{\noindent\hfill
%\begin{minipage}[!h]{12cm}
%\singlespacing\fontsize{10}{10}\selectfont
%}{
%\end{minipage}\vspace{.5cm}
%}

%\newcommand{\opt}[5]{\vspace{.5cm} % .......................................................... Cinco opções
%\textbf{a)} #1 \par
%\textbf{b)} #2 \par
%\textbf{c)} #3 \par
%\textbf{d)} #4 \par
%\textbf{e)} #5 \par
%}
% ===========================================================================================================================
\usepackage{helvet}% pacote da fonte helvica
\usepackage{lastpage} % pacote pra pegar nº da última página
\usepackage{adjustbox}% permite ajustar caixas
\usepackage{tabularx} % tabelas diferenciadas. Documentação em:  https://texdoc.org/serve/tabularx.pdf/0
\usepackage{float}% permite maior controle sobre as posições de tableas e figuras
\usepackage{array}% permite controle de larguras mais sofisticado
\usepackage{datetime}
\usepackage{lastpage} % pacote pra pegar nº da última página
\usepackage{ragged2e}% permite outras formatações de texto como o justificado por exemplo
\usepackage[font=sf,labelfont=bf]{caption}% edita as captions das fotos e tabelas. Documentação presente em : https://ctan.math.utah.edu/ctan/tex-archive/macros/latex/contrib/caption/caption.pdf
\DeclareCaptionLabelFormat{andtable}{#1~#2  \&  \tablename~\thetable}% formatação de caption pra figura e tabela juntos

\newcolumntype{g}{>{\columncolor[HTML]{92D050}} c }%m{0.25\textwidth}
% nova formatação de coluna, dessa vez ela está colorida

\usepackage{listings}% pacote para exibição de códigos em espaços ambientados para a linguagem que vc definir. Para mais informações acesse: https://www.overleaf.com/learn/latex/Code_listing
\usepackage{caption}

\usepackage{xcolor} % to access the named colour LightGray
\definecolor{LightGray}{gray}{0.9}

\definecolor{codegreen}{rgb}{0,0.6,0}
\definecolor{codegray}{rgb}{0.5,0.5,0.5}
\definecolor{codepurple}{rgb}{0.58,0,0.82}
\definecolor{backcolour}{rgb}{0.95,0.95,0.92}

%Code listing style named "mystyle"
\lstdefinestyle{mystyle}{
  backgroundcolor=\color{backcolour},   commentstyle=\color{codegreen},
  keywordstyle=\color{magenta},
  numberstyle=\tiny\color{codegray},
  stringstyle=\color{codepurple},
  basicstyle=\ttfamily\footnotesize,
  breakatwhitespace=false,         
  breaklines=true,                 
  captionpos=b,                    
  keepspaces=true,                 
  numbers=left,                    
  numbersep=5pt,                  
  showspaces=false,                
  showstringspaces=false,
  showtabs=false,                  
  tabsize=2
} % Para saber mais coisas sobre os estilos para a exibição de código veja mais em : https://www.overleaf.com/learn/latex/Code_listing#Reference_guide

%"mystyle" code listing set
\lstset{style=mystyle}

\renewcommand\lstlistingname{Algoritmo}
%\usepackage[
%backend=biber,
%stylename=ieee,
%style=numeric,
%sorting=ynt
%]{biblatex}



%\setkomafont{caption}{\small\fontfamily{phv}\selectfont}
%\setkomafont{captionlabel}{\usekomafont{caption}}
% Novos contadores ==========================================================================================================
%\newcounter{quest} % ......................................................................... Nova questão
%\setcounter{quest}{1}


%\newenvironment{quest}{
%\noindent\textbf{Questão \thequest)\stepcounter{quest}}
%}{
%\vspace{.5cm}
%}
% ===========================================================================================================================

%\smallskip
%Adds a 3pt space plus or minus 1pt depending on other factors (document type, available space, etc)
%\medskip
%Adds a 6pt space plus or minus 2pt depending on other factors (document type, available space, etc)
%\bigskip
%Adds a 12pt space plus or minus 4pt depending on other factors (document type, available space, etc)

\newcounter{rfreq}
\newenvironment{rfreq}[1]
{
    \refstepcounter{rfreq}
    \subsubsection{[RF\ifnum\value{rfreq}<10 0\fi\arabic{rfreq}] \textcolor{black}{[#1]}}
    \begin{flushright}
    \begin{minipage}{0.9\textwidth}
}
{
    \end{minipage}
    \end{flushright}
}

\newcounter{nfreq}
\newenvironment{nfreq}[1]
{
    \refstepcounter{nfreq}
    \subsubsection{[NF\ifnum\value{nfreq}<10 0\fi\arabic{nfreq}] \textcolor{black}{[#1]}}
    \begin{flushright}
    \begin{minipage}{0.9\textwidth}
}
{
    \end{minipage}
    \end{flushright}
}

\usepackage{enumitem} %Permite enumerar com letras
\usepackage{longtable} %Tabelas continuas por múltiplas páginas
\usepackage{mwe}
